\subsection{Properties of p-th power integrable functions}

THEOREM 4. --- Let F and G be two Banach spaces, u a continuous linear mapping of F into G. For every function \(f \in L_{F}^{p}\) , the composite function \(u \circ f\) belongs to \(L_{G}^{p}\) .

Let \(\mathbf{f} \in \mathcal{L}_{\mathrm{F}}^{p}\) ; for every \(\varepsilon > 0\) , there exists a function \(\mathbf{g} \in \mathcal{K}_{\mathrm{F}}\) such that \(\mathrm{N}_p(\mathbf{f} - \mathbf{g}) \leqslant \varepsilon\) ; since \(|\mathbf{u} \circ \mathbf{f} - \mathbf{u} \circ \mathbf{g}| \leqslant \| \mathbf{u} \| \cdot |\mathbf{f} - \mathbf{g}|\) , we have

\[
\mathrm{N} _ {p} (\mathbf {u} \circ \mathbf {f} - \mathbf {u} \circ \mathbf {g}) \leqslant \| \mathbf {u} \| \cdot \mathrm{N} _ {p} (\mathbf {f} - \mathbf {g}) \leqslant \varepsilon \| \mathbf {u} \|,
\]

and since \(u \circ g\) is continuous with compact support, the theorem is proved.

COROLLARY 1. --- Let \(\mathbf{a}'\) be a continuous linear form on \(\mathbf{F}\) ; if \(\mathbf{f} \in \mathcal{L}_{\mathrm{F}}^{p}\) , the numerical function \(x \mapsto \langle \mathbf{f}(x), \mathbf{a}' \rangle\) (denoted by \(\langle \mathbf{f}, \mathbf{a}' \rangle\) ) belongs to \(\mathcal{L}^p\) .

COROLLARY 2. --- Given \(n\) points \(\mathbf{a}_k\) of \(F (1 \leqslant k \leqslant n)\) , and \(n\) numerical functions \(f_k (1 \leqslant k \leqslant n)\) belonging to \(\mathcal{L}^p\) , the function \(\mathbf{f} = \sum_{k=1}^{n} \mathbf{a}_k f_k\) belongs to \(\mathcal{L}_F^p\) .

This follows from the fact that the mapping \(t \mapsto at\) of R into F is continuous.

PROPOSITION 9. --- Let F be an n-dimensional vector space over R, and let \((\mathbf{e}_{k})_{1\leqslant k\leqslant n}\) be a basis of F. For a function \(f=\sum_{k=1}^{n}\mathbf{e}_{k}f_{k}\) to belong to \(L_{F}^{p}\) , it is necessary and sufficient that each of the numerical functions \(f_{k}\) belong to \(L^{p}\) .

This follows at once from Cors. 1 and 2 of Th. 4.

PROPOSITION 10. --- In the space \(\mathcal{L}_{\mathrm{F}}^{p}\) , the linear subspace formed by the (finite) linear combinations \(\sum_{k} \mathbf{a}_{k} f_{k}\) , where \(\mathbf{a}_{k} \in \mathbf{F}\) and the \(f_{k}\) are continuous numerical functions with compact support, is dense (for the topology of convergence in mean of order \(p\) ).

The set \(H_{F}\) of continuous mappings of X into F with compact support is by definition dense in \(L_{F}^{p}\) . On the other hand, every function \(g \in H_{F}\) may be uniformly approximated by functions of the form \(\sum_{k} a_{k} f_{k}\) , where the \(f_{k}\) are continuous functions with support contained in a fixed compact neighborhood of the support of g (Ch. III, §1, No.~2, Lemma 2); it follows (No.~3, Prop. 4) that g is in the closure in \(L_{F}^{p}\) of the set of \(\sum_{k} a_{k} f_{k}\) , whence the proposition.

PROPOSITION 11. --- If a function f belongs to \(L_{F}^{p}\) , then the function \(|f|\) belongs to \(L^{p}\) , and the mapping \(f \mapsto |f|\) of \(L_{F}^{p}\) into \(L^{p}\) is uniformly continuous (for the topology of convergence in mean of order p).

For every \(\varepsilon > 0\) there exists a continuous function \(\mathbf{g}\) with compact support, such that \(\mathrm{N}_p(\mathbf{f} - \mathbf{g}) \leqslant \varepsilon\) ; since \(||\mathbf{f}| - |\mathbf{g}|| \leqslant |\mathbf{f} - \mathbf{g}|\) , we have \(\mathrm{N}_p(|\mathbf{f}| - |\mathbf{g}|) \leqslant \varepsilon\) , which proves that \(|\mathbf{f}| \in \mathcal{L}^p\) . On the other hand, if \(\mathbf{f}_1, \mathbf{f}_2\) are two functions in \(\mathcal{L}_{\mathrm{F}}^p\) then \(\mathrm{N}_p(|\mathbf{f}_1| - |\mathbf{f}_2|) \leqslant \mathrm{N}_p(\mathbf{f}_1 - \mathbf{f}_2)\) , which shows that \(\mathbf{f} \mapsto |\mathbf{f}|\) is a uniformly continuous mapping.

PROPOSITION 12. --- For a numerical function f to belong to \(L^{p}\) , it is necessary and sufficient that each of the functions \(f^{+}\) and \(f^{-}\) belong to \(L^{p}\) .

The condition is sufficient since \(f = f^{+} - f^{-}\) ; it is necessary, because if \(f \in \mathcal{L}^p\) then \(|f| \in \mathcal{L}^p\) (Prop. 11).

COROLLARY. --- The upper (resp. lower) envelope of a finite family of functions in \(\mathcal{L}^{p}\) belongs to \(\mathcal{L}^{p}\) .

